Como mencionamos en la sección anterior, $ASV$ utiliza el grafo causal asociado al problema para definir una función $w$, que nos va a permitir filtrar permutaciones no deseadas que no respeten la causalidad. Luego, cuando tenemos una permutación $\topo$ que sí nos interesa, vamos a querer realizar la operación $\charactheristicFunction(\pi_{<i})$, la cual consiste en evaluar la esperanza teniendo en cuenta nuestra distribución elegida y un valor fijo para los features en $\pi_{<i}$. Para realizar este cálculo necesitamos entender cómo se calculan probabilidades en una red bayesiana y cuál es su complejidad. 

\subsection{Redes Bayesianas}
Una Red Bayesiana $\aBayesianNetwork$ es un $DAG$ donde cada nodo representa una variable y las aristas representan dependencias condicionales entre ellas. Por ejemplo, un arco $A \rightarrow B$ indica que la variable $B$ depende\footnote{Aunque no necesariamente implica una relación causal entre las variables, ya que la red puede no coincidir con el DAG causal subyacente.} de $A$, lo cual se modela mediante la probabilidad condicional $P(B \mid A)$. Un nodo puede tener múltiples padres, por lo que su distribución se especifica mediante una \emph{Tabla de Probabilidad Condicional} (CPT), que define la probabilidad de cada uno de sus valores dado el valor de sus padres. Formalmente, una Red Bayesiana sobre un conjunto de variables $X$ es una tupla $\aBayesianNetwork = (X, E, Pr)$, donde $(X,E)$ es un DAG y $Pr$ define para cada variable $x \in X$ su distribución condicional $Pr(x \mid \parents(x))$. Cada variable posee un conjunto finito de estados posibles y, si tiene padres $B_1,\dots,B_n$, su comportamiento probabilístico queda definido por una CPT $Pr(x \mid B_1,\dots,B_n)$.

Una propiedad fundamental de las redes bayesianas es que permiten descomponer la distribución de probabilidad conjunta utilizando dependencias locales. Si $\aBayesianNetwork$ está definida sobre el conjunto de variables $U=\{A_1,\dots,A_n\}$, entonces la distribución conjunta puede expresarse como
\[
Pr(U) = \prod_{i=1}^{n} Pr(A_i \mid \parents(A_i)).
\]
Esta factorización simplifica el cálculo de probabilidades y permite realizar inferencia probabilística en la red, es decir, calcular distribuciones condicionadas a evidencia sobre algunas variables. La complejidad de este proceso depende de la estructura del grafo: en redes con forma de \emph{polytree} (DAG cuyo grafo subyacente es un árbol) la inferencia puede realizarse en tiempo polinomial \cite{pearl1986bayesianInference}, mientras que en redes generales es $\sharpPhard$. Un algoritmo común para realizar inferencia marginal es \emph{Variable Elimination} \cite{variableElimination}, cuya complejidad depende del \emph{treewidth} de la red; en redes con treewidth acotado, la inferencia es polinomial en el tamaño de la red. Por esta razón, en este trabajo nos enfocamos en redes bayesianas con estructura de polytree.



\subsection{Predicción promedio en árboles de decisión}

Los árboles de decisión son modelos ampliamente utilizados en Inteligencia Artificial por su capacidad para realizar predicciones y clasificaciones a partir de los features de entrada. Su estructura jerárquica permite representar decisiones mediante una secuencia de preguntas simples, lo que contribuye a la interpretabilidad del modelo, aunque el camino recorrido no siempre constituye la mejor explicación de la predicción obtenida.\footnote{Aunque no siempre el camino de un árbol es la mejor explicación, ya que puede haber features redundantes en el camino que no sean claves para la predicción realizada. \cite{audemard2021explanatorypowerdecisiontrees}} En este trabajo nos interesa calcular la predicción promedio de un árbol de decisión cuando la distribución de probabilidad sobre los features está dada por una red bayesiana, con el objetivo de estudiar la tratabilidad del cálculo de ASV. Para ello, consideramos primero árboles de decisión binarios con features binarios. Para calcular la predicción promedio proponemos un algoritmo que recorre recursivamente todas las ramas del árbol acumulando en \texttt{pathCondition} las asignaciones inducidas por las decisiones tomadas en cada nodo. Al llegar a una hoja, calcula la probabilidad de haber alcanzado esa hoja dada la evidencia utilizando la red bayesiana, y multiplica dicho valor por la predicción almacenada en la hoja. La predicción promedio se obtiene entonces como la suma de los aportes de todas las hojas compatibles con la evidencia.

La complejidad del algoritmo es $O(i|V| + l \cdot varElim)$, donde $i$ es la cantidad de nodos internos, $l$ la cantidad de hojas y $varElim$ representa el costo del algoritmo de inferencia por eliminación de variables. En el caso de los polytrees, esta inferencia es polinomial, por lo que el cálculo de la predicción promedio también lo es en función del tamaño del árbol de decisión. Además, este enfoque permite capturar el efecto de evidencia sobre variables que no aparecen explícitamente en el árbol, pero que modifican la inferencia en la red bayesiana y, en consecuencia, la predicción promedio. Esta idea puede extenderse a features no binarios, aunque en la implementación actual la inferencia deja de ser polinomial, ya que cada condición de camino puede involucrar conjuntos de valores y su probabilidad debe calcularse mediante sumas de múltiples consultas puntuales. Si bien sería posible optimizar este procedimiento modificando el algoritmo de \textit{Variable Elimination} o incorporando nodos intermedios que representen la evidencia, no profundizamos en esa dirección porque excede el objetivo principal de esta tesis.